\subsection{Approximation to convex functions}

PROPOSITION 5. --- Let X be a closed convex set in a locally convex space E. Then every lower semi-continuous convex function f defined in X is the upper envelope of a family of functions that are the restrictions to X of continuous affine linear functions in E.

For, the set \(A \subset E \times R\) of points \((x, t)\) such that \(x \in X\) and \(t \geqslant f(x)\) is convex (II, p.~17, prop. 19) and closed, since the function \((x, t) \mapsto f(x) - t\) is lower semicontinuous. Then let x be any point of X and let \(a \in R\) be such that \(a < f(x)\) . By cor. 1 of II, p.~38, there exists a closed hyperplane H in \(E \times R\) , that contains \((x, a)\) and does not meet A. Every linear continuous form on \(E \times R\) being of the form

\[
(z, t) \rightarrow u (z) + \lambda t,
\]

where \(\lambda\in R\) and u is a continuous linear form on E, it follows that H has an equation of the form \(u(z)+\lambda t=\alpha\) , and as H contains \((x,a)\) we have \(\alpha=u(x)+\lambda a\) . Now the point \((x,f(x))\in\mathbf{A}\) does not belong to H and therefore \(\lambda\neq0\) . Dividing by \(-\lambda\) , if necessary, we can write the equation of H as \(t-a=u(z-x)\) . As \(f(x)-a>0\) , we have, therefore, \(f(z)>u(z-x)+a\) for all \(z\in X\) and this proves the proposition.

Remarks. --- 1) It follows from prop. 5 that f is the upper envelope of a directed increasing family of functions that are the restrictions to X of functions which are continuous and convex in E.

\begin{enumerate}
\def\labelenumi{\arabic{enumi})}
\setcounter{enumi}{1}
\tightlist
\item
  Suppose further that X is a closed convex cone with vertex 0 and that f is positively homogeneous. Then f is the upper envelope of a family of functions which are the restrictions to X of continuous linear forms in E. For, let \((u_{\alpha})\) be a family of continuous affine linear functions in E of which the restrictions to X have f as their upper envelope. Put \(u_{\alpha} = v_{\alpha} + \lambda_{\alpha}\) , where \(\lambda_{\alpha} \in R\) , and where \(v_{\alpha}\) is a continuous linear form in E. We have \(\lambda_{\alpha} = u_{\alpha}(0) \leqslant f(0) = 0\) . On the other hand, if \(x \in X\) , we have for every \(\mu > 0\) ,
\end{enumerate}

\[
\mu^ {- 1} \lambda_ {\alpha} + v _ {\alpha} (x) = \mu^ {- 1} (\lambda_ {\alpha} + v _ {\alpha} (\mu x)) = \mu^ {- 1} u _ {\alpha} (\mu x) \leqslant \mu^ {- 1} f (\mu x) = f (x)
\]

therefore \(u_{\alpha} \leqslant v_{\alpha} \leqslant f\) in X so that \(f\) is the upper envelope of the \(v_{\alpha}\) .

\begin{enumerate}
\def\labelenumi{\arabic{enumi})}
\setcounter{enumi}{2}
\tightlist
\item
  The restriction to X of a continuous affine function in E is a function that is affine in X (i.e.~both concave and convex II, p.~17); but it may be the case that there exist continuous affine functions in a compact convex set \(\mathbf{X} \subset \mathbf{E}\) , that are not the restrictions to X of continuous affine functions in E (II, p.~78, exerc. II, c)). However:
\end{enumerate}

PROPOSITION 6. --- Let f be an upper semi-continuous affine function in a compact convex set X, of the Hausdorff locally convex space E. Let L be the set of restrictions to X of continuous affine functions in E; the set L' of the \(h \in L\) such that \(h(x) > f(x)\) for all \(x \in X\) , is then decreasing directed and its lower envelope is equal to f.

We may suppose that X is non-empty. Let u, v be two elements of L, such that \(u(x) > f(x)\) and \(v(x) > f(x)\) for all \(x \in X\) , and let b be a constant that is an upper bound of u and v. Let U (resp. V) be the compact convex set of points \((x, t)\) of \(X \times R\) such that \(u(x) \leqslant t \leqslant b\) (resp. \(v(x) \leqslant t \leqslant b\) ), and let F be the set of \((x, t) \in X \times R\) such that \(t \leqslant f(x)\) ; F is convex and closed in \(X \times R\) . The convex envelope K of \(U \cup V\) does not meet F, since \(U \cup V\) is contained in the set of \((x, t) \in X \times R\) such that \(f(x) < t\) , a set which is convex and does not meet F. As K is compact (II, p.~14, prop. 15), we can separate F strictly from K by a closed hyperplane H in \(E \times R\) . For every \(x \in X\) , the hyperplane H separates \((x, f(x))\) strictly from \((x, b)\) , and therefore meets the line \(\{x\} \times R\) in a single point \(w(x)\) ; thus H is the graph of a continuous affine function whose restriction w to X is a member of L, that is a lower bound for u and v and that satisfies the inequality \(w(x) > f(x)\) for all \(x \in X\) . This proves that the set \(L'\) is decreasing directed. Prop. 5 of II, p.~39, applied to -f shows that f is the lower envelope of \(L'\) .

COROLLARY. --- Let \(f\) be a continuous affine function in \(\mathbf{X}\) ; then there exists a sequence \((h_n)\) of elements of \(\mathbf{L}\) which converges uniformly to \(f\) in \(\mathbf{X}\) .

For, prop. 6 and Dini's theorem (GT, X, § 4.1, th. 1) show that for all \(n\) there exists \(h_n \in \mathbf{L}\) such that \(f \leqslant h_n \leqslant f + 1/n\) .

